\section{Discussion}

We have seen that the graphs arising from the Penrose dodecahedron share many unusual 
properties which hint to the fact that they might be quantum isomorphic,
and moreover, no property was found that definitely rules out this conjectured quantum isomorphism.

Below we attempt to apply the method from~\cite{schmidt2022quantumisomorphicstronglyregular}
in order to directly construct this quantum isomorphism,
and show that unfortunately, there are several structural hurdles
that complicates the search for a suitable magic unitary.

\subsection{Applying the Approach from~\cite{schmidt2022quantumisomorphicstronglyregular}}
We firstly observe that $Q^*$ has independence number $10$
and, by straightforward computation, see that there are a total of $2880$ maximally independent sets,
which translates to the fact that there are $2880$ possible partitions of $V(Q)$ into 10 4-cliques.
However, no subgroup of $\operatorname{Aut}(Q)$ exists that induces 10 orbits comprising of such cliques,
unlike was found for $G_{E_8}$ in~\cite{schmidt2022quantumisomorphicstronglyregular}.

Secondly, it is impossible to apply the construction of $G^\omega$ in~\cite{schmidt2022quantumisomorphicstronglyregular}
to $Q$ in order to receive a graph isomorphic to $Q^*$.
This is due to the following observations.
Firstly, the induced subgraph of two disjoint cliques in $Q$ do not share the property of the analogous subgraphs of $G_{E_8}$ that for each vertex inside clique $V_i$, there are exactly as many neighbours in $V_j$ as there are non-neighbours. 
This allowed to switch the adjacency relations while preserving the edge count.
However, for both $Q$ and $Q^*$, each vertex in one clique $V_i$ has exactly one neighbour and three non-neighbours in each $V_j$,
meaning that applying this form of switching changes the edge count.
However, this contradicts the fact that $Q$ and $Q^*$, which have the same edge count, are related via this form of switching.

More importantly, the duality between $Q$ and $Q^*$ that vertices in one graph are 4-cliques in the other 
forbids us to fix a partition of 10 4-cliques, since any graph resulting from this procedure would still comprise of 10 disjoint 4-cliques. 
Assuming that such a graph is isomorphic to $Q^*$ implies that $Q$ has independence number 10, which is a contradiction.
%
Therefore, the magic unitary corresponding to the potential quantum isomorphism cannot be a block diagonal matrix
with 10 blocks of size 4 by 4.

Moreover, it is unresolved whether either of the graphs has quantum symmetry,
as it was shown by exhaustive enumeration that there are no disjoint automorphisms in $\operatorname{Aut}(Q)$,
which would imply the existence of a quantum automorphism, as proved in~\cite{Schmidt_2020}.